\section{The Model: Estimates, Fit and Decompositions}\label{app:modelfit}

This appendix collects the tables behind Section~\ref{sec:model}: the fit of every model of the lattice
(Table~\ref{tab:modellattice}), the estimates of the selected model (Table~\ref{tab:modelest}), its fit to all 66
targeted moments (Table~\ref{tab:modelfit}), the pay-off of parental income by age and parental tercile
(Table~\ref{tab:modelcfh4}) and the sensitivity of the race parameters to the moments (Table~\ref{tab:modelsens}). All
are generated from the estimation output by \texttt{model/tables.py}.

\begin{table}[!htbp]
\centering
\caption{The hypothesis lattice: fit of every combination of race parameters}
\label{tab:modellattice}
\resizebox{\ifdim\width>\textwidth\textwidth\else\width\fi}{!}{%
\begin{tabular}{lccccccccc}
\toprule
Model & H1 & H2 & H3 & H4 & $Q$ & $Q$ over ten & W $-$ B IGE gap & W $-$ B IGE gap & $Q$ of the \\
 & & & & & & seeds: mean (s.d.) & 30--34 & 45--49 & 12 IGE moments \\
\midrule
\multicolumn{10}{l}{\textit{Without H4}} \\
A without H4 &  &  &  &  & 182.0 & 184.9 (2.5) & 0.079 & 0.142 & 36.9 \\
E without H4 &  & $\checkmark$ & $\checkmark$ &  & 160.6 & 163.7 (3.5) & 0.079 & 0.154 & 32.7 \\
\midrule \multicolumn{10}{l}{\textit{With H4}} \\
A &  &  &  & $\checkmark$ & 138.4 & 137.8 (2.7) & 0.174 & 0.275 & 7.3 \\
B & $\checkmark$ &  &  & $\checkmark$ & 135.0 & 134.7 (2.1) & 0.169 & 0.271 & 6.8 \\
F &  & $\checkmark$ &  & $\checkmark$ & 136.4 & 136.1 (2.4) & 0.169 & 0.265 & 7.3 \\
G &  &  & $\checkmark$ & $\checkmark$ & 138.2 & 137.6 (2.4) & 0.174 & 0.276 & 7.0 \\
C & $\checkmark$ & $\checkmark$ &  & $\checkmark$ & 133.8 & 133.7 (1.9) & 0.168 & 0.266 & 6.8 \\
H & $\checkmark$ &  & $\checkmark$ & $\checkmark$ & 135.0 & 134.7 (2.1) & 0.169 & 0.271 & 6.8 \\
E &  & $\checkmark$ & $\checkmark$ & $\checkmark$ & 127.6 & 127.6 (2.3) & 0.169 & 0.282 & 4.9 \\
D & $\checkmark$ & $\checkmark$ & $\checkmark$ & $\checkmark$ & 127.1 & 126.2 (2.2) & 0.167 & 0.280 & 5.0 \\
\midrule
PSID & & & & & & & 0.153 & 0.330 & \\
\bottomrule
\end{tabular}}
\par\smallskip
\begin{minipage}{\textwidth}\footnotesize \textit{Notes:} Each model is estimated by SMM on the 66 moments with 20 simulated sons per PSID son (12,000 CMA-ES evaluations, each started from the best solution of the models nested in it). $Q$: the criterion at the estimation draws, and its mean (standard deviation) over ten simulation seeds. IGE gaps and the $Q$ of the IGE moments at the estimation draws. IGE gaps: White $-$ Black coefficient of the son's log wage on log parental income, simulated with the PSID observation scheme. H1 access ($\alpha_3$), H2 pay-off of schooling in accumulation, H3 entry return to schooling, H4 pay-off of parental income in accumulation. Model H's solution is model B's (see the note to Table~\ref{tab:modelhyp}). Code: \texttt{model/hpc\_ige\_chain.sh}, \texttt{model/lattice\_q.py}.\end{minipage}
\end{table}

\begin{table}[!htbp]
\centering
\caption{Estimates of the selected model (E: H2, H3 and H4; job finding out of work by race)}
\label{tab:modelest}
\resizebox{\ifdim\width>\textwidth\textwidth\else\width\fi}{!}{%
\begin{tabular}{lc}
\toprule
Parameter & Estimate (s.e.) \\
\midrule
\multicolumn{2}{l}{\textit{Race}} \\
\quad $\mu_B$, racial wage level & $-0.1458^{***}$ (0.0248) \\
\quad $\gamma_1(B) - \gamma_1(W)$, pay-off of schooling (H2) & $-0.0080^{***}$ (0.0029) \\
\quad $b_1(B) - b_1(W)$, entry return to schooling (H3) & $0.0341^{**}$ (0.0136) \\
\quad $\gamma_x(B) - \gamma_x(W)$, pay-off of parental income (H4) & $-0.0194^{***}$ (0.0054) \\
\multicolumn{2}{l}{\textit{Parental income}} \\
\quad $b_x$, starting level & $0.0938^{***}$ (0.0332) \\
\quad $\alpha_1$, access to $G$ jobs & $0.2073^{*}$ (0.1186) \\
\quad $\gamma_x(W)$, accumulation & $0.0186^{***}$ (0.0072) \\
\quad $\delta_1$, displacement & $-0.2033^{***}$ (0.0554) \\
\multicolumn{2}{l}{\textit{Schooling}} \\
\quad $b_1(W)$, starting level & $0.0764^{***}$ (0.0117) \\
\quad $\alpha_2$, access to $G$ jobs & $0.6862^{***}$ (0.0455) \\
\quad $\gamma_1(W)$, accumulation & $0.0125^{***}$ (0.0024) \\
\multicolumn{2}{l}{\textit{Jobs, search frictions and job loss}} \\
\quad $\alpha_0$, access to $G$ jobs & $-5.5035^{***}$ (0.2147) \\
\quad $\gamma_O$, accumulation on $O$ jobs & $0.0878^{***}$ (0.0114) \\
\quad $\gamma_G$, accumulation on $G$ jobs & $0.1629^{***}$ (0.0351) \\
\quad $\ln\rho$, decline of accumulation with age & $-2.3960$ (0.1014) \\
\quad $\ln\sigma_m$, match component & $-2.1313$ (0.1968) \\
\quad logit $\lambda_e$ at 30, on-the-job search friction & $0.1579$ (0.1688) \\
\quad $\lambda_a$, its age slope (per decade) & $-0.7308$ (0.1161) \\
\quad logit $\bar\lambda_u$, job finding out of work & $0.3099$ (0.0703) \\
\quad $\Delta\ell_u$, job finding out of work, Black $-$ White (logit) & $-0.5389^{***}$ (0.0946) \\
\quad $\delta_{0B}$, displacement, Black & $-2.3353^{***}$ (0.0495) \\
\quad $\delta_{0W}$, displacement, White & $-2.8068^{***}$ (0.0475) \\
\quad $\delta_a$, its age slope (per decade) & $-0.4219^{***}$ (0.0416) \\
\multicolumn{2}{l}{\textit{Starting level and shocks}} \\
\quad $b_0$, starting level & $1.9033^{***}$ (0.0522) \\
\quad $\ln\sigma_0$, starting level & $-0.8759$ (0.0267) \\
\quad $\ln\sigma_\eta$, permanent shock & $-2.6922$ (0.0764) \\
\quad $\ln\sigma_\varepsilon$, transitory shock & $-1.5213$ (0.0292) \\
\midrule
$Q$ (66 moments), estimation draws & 95.9 \\
\bottomrule
\end{tabular}}
\par\smallskip
\begin{minipage}{\textwidth}\footnotesize \textit{Notes:} Simulated method of moments, 66 moments, diagonal weights equal to the inverse cluster-bootstrap variances (500 draws of 1968 families), 20 simulated sons per PSID son. Sandwich standard errors $(1 + 1/R)\,(G'WG)^{-1}G'W\Omega WG(G'WG)^{-1}$ with $\Omega$ the bootstrap covariance of the moments; $^{***}$ $p<0.01$, $^{**}$ $p<0.05$, $^{*}$ $p<0.1$ (not shown for parameters on a log or logit scale). Accumulation parameters per year; parental income and schooling are centred at the sample mean and at 12 years. Fixed: $\beta = 0.95$, $b = 0.4$, cyclical elasticities $\mp 0.10$, $\alpha_3 = 0$ (H1 rejected), the match premium of $G$ jobs $\mu_G = 0$ and depreciation out of work $d_u = 0$ (Section~\ref{sec:model}). Code: \texttt{model/estimate.py}; output \texttt{model/output/final\_x4/estimates\_E.json}.\end{minipage}
\end{table}

\begin{table}[!htbp]
\centering
\caption{Fit of the selected model (E): targeted moments}
\label{tab:modelfit}
\resizebox{\ifdim\width>\textwidth\textwidth\else\width\fi}{!}{%
\begin{tabular}{lcccc}
\toprule
Moment & PSID & (s.e.) & Model & $t$ \\
\midrule
\multicolumn{5}{l}{\textit{Persistence by age (log-wage IGE)}} \\
\quad IGE, 25--29, White & $0.205$ & (0.017) & $0.223$ & $1.1$ \\
\quad IGE, 30--34, White & $0.335$ & (0.023) & $0.322$ & $-0.6$ \\
\quad IGE, 35--39, White & $0.381$ & (0.030) & $0.386$ & $0.2$ \\
\quad IGE, 40--44, White & $0.439$ & (0.036) & $0.428$ & $-0.3$ \\
\quad IGE, 45--49, White & $0.470$ & (0.049) & $0.457$ & $-0.3$ \\
\quad IGE, 50--55, White & $0.450$ & (0.054) & $0.483$ & $0.6$ \\
\quad IGE, 25--29, Black & $0.148$ & (0.021) & $0.140$ & $-0.4$ \\
\quad IGE, 30--34, Black & $0.182$ & (0.024) & $0.154$ & $-1.2$ \\
\quad IGE, 35--39, Black & $0.182$ & (0.030) & $0.172$ & $-0.4$ \\
\quad IGE, 40--44, Black & $0.146$ & (0.034) & $0.169$ & $0.7$ \\
\quad IGE, 45--49, Black & $0.140$ & (0.044) & $0.176$ & $0.8$ \\
\quad IGE, 50--55, Black & $0.151$ & (0.059) & $0.173$ & $0.4$ \\
\multicolumn{5}{l}{\textit{Return to schooling}} \\
\quad return to schooling, 25--29, White & $0.047$ & (0.005) & $0.046$ & $-0.2$ \\
\quad return to schooling, 45--49, White & $0.125$ & (0.010) & $0.138$ & $1.2$ \\
\quad return to schooling, 25--29, Black & $0.067$ & (0.007) & $0.073$ & $1.0$ \\
\quad return to schooling, 45--49, Black & $0.119$ & (0.012) & $0.110$ & $-0.7$ \\
\multicolumn{5}{l}{\textit{Careers}} \\
\quad share in $G$ jobs at 25--34, bottom tercile, Black & $0.031$ & (0.005) & $0.032$ & $0.2$ \\
\quad share in $G$ jobs at 25--34, top tercile, Black & $0.117$ & (0.022) & $0.100$ & $-0.8$ \\
\quad share in $G$ jobs at 25--34, bottom tercile, White & $0.066$ & (0.013) & $0.047$ & $-1.5$ \\
\quad share in $G$ jobs at 25--34, top tercile, White & $0.180$ & (0.011) & $0.181$ & $0.1$ \\
\quad same-employer growth, 25--34, $O$ jobs & $0.055$ & (0.003) & $0.054$ & $-0.4$ \\
\quad same-employer growth, 25--34, $G$ jobs & $0.075$ & (0.007) & $0.078$ & $0.4$ \\
\quad displacement rate, 25--34, Black & $0.173$ & (0.009) & $0.170$ & $-0.4$ \\
\quad displacement rate, 25--34, White & $0.102$ & (0.005) & $0.107$ & $0.9$ \\
\quad displacement rate, 35--55, Black & $0.111$ & (0.006) & $0.115$ & $0.6$ \\
\quad displacement rate, 35--55, White & $0.076$ & (0.004) & $0.074$ & $-0.6$ \\
\quad move rate, 25--34, White & $0.131$ & (0.004) & $0.132$ & $0.3$ \\
\quad move rate, 35--55, White & $0.091$ & (0.004) & $0.087$ & $-1.1$ \\
\quad wage growth at a move, 25--34, Black & $-0.006$ & (0.019) & $0.060$ & $3.5$ \\
\quad wage growth at a move, 25--34, White & $0.052$ & (0.010) & $0.071$ & $2.0$ \\
\quad zero-hours years, Black & $0.109$ & (0.006) & $0.109$ & $-0.1$ \\
\quad zero-hours years, White & $0.039$ & (0.003) & $0.039$ & $0.2$ \\
\quad variance of log wages, 45--49, Black & $0.345$ & (0.021) & $0.417$ & $3.5$ \\
\quad variance of log wages, 45--49, White & $0.503$ & (0.023) & $0.470$ & $-1.5$ \\
\quad variance of stayers' growth & $0.121$ & (0.005) & $0.103$ & $-3.3$ \\
\quad autocovariance of stayers' growth & $-0.037$ & (0.003) & $-0.048$ & $-3.6$ \\
\bottomrule
\end{tabular}}
\par\smallskip
\begin{minipage}{\textwidth}\footnotesize \textit{Notes:} PSID sons (Black and White male heads, 25--55), cluster-bootstrap standard errors; model moments from the simulated panel with the PSID's observation scheme at the estimate (20 simulated sons per PSID son). $t$ = (model $-$ PSID) / s.e. Rates per interview pair, growth annualised. All 66 moments: \texttt{model/output/final\_x4/estimates\_E.json}.\end{minipage}
\end{table}

\begin{table}[!htbp]
\centering
\caption{The pay-off of parental income in accumulation: by age and parental income (model E)}
\label{tab:modelcfh4}
\resizebox{\ifdim\width>\textwidth\textwidth\else\width\fi}{!}{%
\begin{tabular}{lcccccc}
\toprule
 & Black IGE & H4 off: & W $-$ B gap & H4 off: & White IGE & No pay-off: \\
 & baseline & change & baseline & change & baseline & change \\
\midrule
\multicolumn{7}{l}{\textit{A. Log-wage IGE by age}} \\
\quad 25--29 & 0.140 & $0.087^{***}$ (0.024) & 0.083 & $-0.087^{***}$ (0.024) & 0.223 & $-0.067^{***}$ (0.025) \\
\quad 30--34 & 0.154 & $0.120^{***}$ (0.032) & 0.168 & $-0.120^{***}$ (0.032) & 0.322 & $-0.104^{***}$ (0.039) \\
\quad 35--39 & 0.172 & $0.141^{***}$ (0.038) & 0.214 & $-0.141^{***}$ (0.038) & 0.386 & $-0.127^{***}$ (0.048) \\
\quad 40--44 & 0.169 & $0.156^{***}$ (0.042) & 0.259 & $-0.156^{***}$ (0.042) & 0.428 & $-0.143^{***}$ (0.053) \\
\quad 45--49 & 0.176 & $0.166^{***}$ (0.045) & 0.281 & $-0.166^{***}$ (0.045) & 0.457 & $-0.155^{***}$ (0.058) \\
\quad 50--55 & 0.173 & $0.173^{***}$ (0.047) & 0.309 & $-0.173^{***}$ (0.047) & 0.483 & $-0.164^{***}$ (0.061) \\
\midrule \multicolumn{7}{l}{\textit{B. Mean log wage by parental-income tercile}} \\
 & Black & H4 off: change & & & White & No pay-off: change \\
\quad 25--34, bottom tercile & 2.335 & $-0.0880^{***}$ (0.0239) & & & 2.458 & $0.0654^{***}$ (0.0247) \\
\quad 25--34, middle tercile & 2.456 & $0.0006^{***}$ (0.0002) & & & 2.668 & $-0.0080^{***}$ (0.0030) \\
\quad 25--34, top tercile & 2.558 & $0.0533^{***}$ (0.0144) & & & 2.845 & $-0.0559^{***}$ (0.0210) \\
\quad 45--49, bottom tercile & 2.642 & $-0.1404^{***}$ (0.0380) & & & 2.774 & $0.1074^{***}$ (0.0402) \\
\quad 45--49, middle tercile & 2.772 & $0.0001$ (0.0002) & & & 3.096 & $-0.0142^{***}$ (0.0053) \\
\quad 45--49, top tercile & 2.915 & $0.0881^{***}$ (0.0239) & & & 3.396 & $-0.1023^{***}$ (0.0383) \\
\bottomrule
\end{tabular}}
\par\smallskip
\begin{minipage}{\textwidth}\footnotesize \textit{Notes:} Model E at its estimates, 20 simulated sons per PSID son, the PSID observation scheme. H4 off: Black sons get White sons' pay-off of parental income in accumulation, $\gamma_x(B) = \gamma_x(W)$. No pay-off: $\gamma_x(B) = \gamma_x(W) = 0$, the race-neutral channel switched off for both races. Parental income is centred at the pooled mean, so a pay-off lowers the wages of sons below the mean and raises those above it. Terciles of parental income over all sons. Delta-method standard errors; $^{***}$ $p<0.01$, $^{**}$ $p<0.05$, $^{*}$ $p<0.1$. Code: \texttt{model/postestimation.py} (\texttt{IGE\_CFS}); output \texttt{model/output/postestimation\_Lx4\_E\_h4.json}.\end{minipage}
\end{table}

\begin{table}[!htbp]
\centering
\caption{Which moments identify the race parameters: sensitivity of the estimates (model E)}
\label{tab:modelsens}
\resizebox{\ifdim\width>\textwidth\textwidth\else\width\fi}{!}{%
\begin{tabular}{lccccccc p{5.2cm}}
\toprule
 & IGE & Returns & Levels, & $G$ & Stayer & Moves, & Zero & Largest loadings \\
 & & & variances & shares & growth & displacement & hours & \\
\midrule
$\gamma_x(B) - \gamma_x(W)$ (H4) & 77 & 10 & 1 & 1 & 9 & 2 & 0 & IGE 30--34 W $-0.21$; IGE 35--39 W $-0.20$; IGE 30--34 B $+0.18$ \\
$\gamma_x(W)$ & 78 & 9 & 1 & 1 & 8 & 3 & 0 & IGE 25--29 B $-0.42$; IGE 25--29 W $-0.30$; IGE 40--44 W $+0.27$ \\
$\gamma_1(B) - \gamma_1(W)$ (H2) & 5 & 63 & 4 & 1 & 25 & 2 & 0 & return 25--29 B $-0.43$; return 25--29 W $+0.29$; return 45--49 B $+0.27$ \\
$b_1(B) - b_1(W)$ (H3) & 11 & 74 & 4 & 0 & 10 & 1 & 0 & return 25--29 B $+0.60$; return 25--29 W $-0.40$; IGE 25--29 B $+0.23$ \\
$\alpha_3$ (H1, model D) & 0 & 0 & 0 & 96 & 0 & 3 & 0 & $G$ share middle tercile B $+0.59$; $G$ share top tercile B $+0.56$; $G$ share middle tercile W $-0.45$ \\
H4 off: change in the gap, 45--49 & 86 & 9 & 1 & 1 & 2 & 1 & 0 & IGE 35--39 W $-0.20$; IGE 40--44 W $-0.19$; IGE 30--34 W $-0.19$ \\
\bottomrule
\end{tabular}}
\par\smallskip
\begin{minipage}{\textwidth}\footnotesize \textit{Notes:} Sensitivity in the sense of \citet{andrews2017measuring}: $\Lambda = (G'WG)^{-1}G'W$, the derivative of the estimate with respect to the data moments, from the numerical Jacobian $G$ at the estimate (20 simulated sons per PSID son) and the diagonal weights $W$. $S_{jk} = \Lambda_{jk}\hat s_k / \mathrm{se}_j$ is the shift of parameter $j$, in its standard errors, when moment $k$ rises by one data standard error. Columns 1--7: the share (\%) of $\sum_k S_{jk}^2$ in each block of moments; last column: the three largest $S_{jk}$ with their signs (race W / B, age band). Last row: the same for the counterfactual change in the W $-$ B IGE gap at 45--49 when Black sons get White sons' pay-off of parental income, from its gradient. Code: \texttt{model/sensitivity.py}; output \texttt{model/output/final\_x4/sensitivity.json}.\end{minipage}
\end{table}


\subsection{Diagnostics of two validation failures}\label{app:modeldiag}

\paragraph{Wage levels after 35.} The model's wages at 35--49
are 0.09--0.15 log points above the PSID for both races ($t$ between 3.4 and 6.5), with the level at 25--29 fitted, and
the failure is growth within sons, not the selection of who has a wage. Among sons with a wage at both 25--29 and
35--39, the PSID's wage grows by 0.332 (0.014) for White and 0.189 (0.018) for Black sons, against 0.439 and 0.303 in
the model; the composition of the two bands moves the PSID's cross-section by 0.03 at most.\footnote{The change of the
mean of son-level band means is split exactly into the change within sons observed in both bands and the two
composition terms. PSID 25--29 to 35--39: White 0.328 = 0.332 $-$ 0.019 + 0.015, Black 0.166 = 0.189 $-$ 0.033 + 0.010
(100 cluster-bootstrap draws); model E: 0.436 and 0.291, of which 0.439 and 0.303 within sons. Code:
\texttt{model/core.py} (\texttt{levels}), \texttt{model/panel\_diag.py}; output
\texttt{model/output/panel\_diag.md}.} Stayers' growth is targeted and fitted at 25--34 and 35--55 ($|t| \le 1.0$);
the wage gain at a move is overstated (0.071 against 0.052 for White and 0.060 against $-0.006$ for Black sons), which
at the PSID's rate of about 0.08 moves a year adds roughly 0.03 (White) and 0.09 (Black) over ten years; the rest of
the excess comes from son-years the pair moments do not cover.

\paragraph{The share in good jobs after 35.} In the PSID the share in $G$ jobs rises with age for White sons (0.144 at 25--29, 0.188 at 30--34, 0.217 at 50--55)
and stays between 0.047 and 0.084 for Black sons; in the model it rises to 30--34 and then falls (White 0.118, 0.141,
0.055; Black 0.050, 0.054, 0.019). Interview pairs show the routes.\footnote{Consecutive interviews one to two years apart from 1985 with a
job type at both, by the son's age at the later one; the change of the $G$ share between the two interviews is split
exactly into the contributions of stayers, movers and displaced sons. Displacement rates by the job type at the
earlier interview. Code: \texttt{model/core.py} (\texttt{routes}), \texttt{model/panel\_diag.py}; output
\texttt{model/output/panel\_diag.md}.} At 35--55 the PSID's share is stable between interviews (net
$+0.004$ a pair for White and $+0.001$ for Black sons): moves and displacement carry no net flow, and the positive part
is reclassification of jobs with the same employer, which the model cannot produce, but which flips back about two
thirds as often as it persists (3.4\% against 5.0\% of three-interview stayer sequences, White), so much of it is coding
noise. In the model the share falls between interviews ($-0.006$ and $-0.002$ a pair) through moves and displacement.
After 35, White sons who move from an $O$ job enter a $G$ job at half the PSID's rate (0.055 against 0.121, s.e.\
0.013), though they do so at the PSID's rate at 25--34 (0.131 against 0.117); and $G$ jobs are lost at half the rate of
$O$ jobs in the PSID (White at 35--55: 0.041 against 0.083; Black 0.045 against 0.113), against 0.79 of it in the model,
where the difference comes only from who holds $G$ jobs. Of the two candidate steps, a profile at the lattice's estimates of
model E picks the second: an access probability that rises with age barely moves the share after 35 (White top
tercile 0.131 to 0.137 at a slope of 0.4 a decade) and lowers it at 25--34, while a separate job-loss rate for $G$ jobs
raises it toward the PSID (to 0.169 for White and 0.102 for Black sons from the top tercile at $-0.6$ on the logit
scale, where $G$ jobs are lost at half the rate of $O$ jobs) and lowers the criterion on the moments extended by the
four shares at 35--55 by 11.7 with every other parameter held at its estimate.\footnote{Code: \texttt{model/core.py}
(switches \texttt{aage} and \texttt{gloss}, moment set \texttt{g35}), data moments
\texttt{model/output/data\_moments\_g35.json}; the joint estimation is pending.} A job-loss rate specific to
good jobs is the next step.
